\documentclass[letterpaper,11pt]{article}
\usepackage[top=0.8in, bottom=0.8in, left=0.9in, right=0.9in]{geometry}
\usepackage[hidelinks]{hyperref}
\usepackage{lmodern}
\usepackage{parskip}
\usepackage{microtype}
\setlength{\parskip}{6pt}
\setlength{\parindent}{0pt}
\begin{document}
\begin{flushleft}
\textbf{\large Ajay Venkatesh} \\
Brooklyn, NY \\
+1 516 585 9013 \\
\href{mailto:av3855@nyu.edu}{av3855@nyu.edu}
\end{flushleft}
\vspace{10pt}
May 21, 2026
\vspace{6pt}

Hiring Manager \\
Diverse Lynx

\vspace{10pt}

Dear Hiring Manager,

My name is Ajay Venkatesh. I'm finishing my M.S. in Computer Engineering at NYU, with production software experience at PwC in Bengaluru, a summer SDE internship at Amazon, and ongoing on-campus work at NYU's Novel AI Technologies. I'm writing about the Java AWS Developer role.

The role's combination, \textbf{Java} backend work delivered on \textbf{AWS} infrastructure, is exactly what I shipped at Amazon. I engineered backend services in \textbf{Java} (Coral, Smithy, Dagger, CBOR) with low-level design patterns and unit / integration testing via \textbf{Mockito}, applied under code-review discipline. I provisioned the supporting infrastructure with \textbf{AWS CDK} (Lambda, API Gateway, DynamoDB, IAM, CloudWatch) and operated multi-stage, multi-region CI/CD with rollback controls. I also built a high-throughput, event-driven processing system on \textbf{AWS ECS Fargate, SQS, SNS} that decoupled distributed message handling under variable load.

On the broader engineering picture: at PwC I owned a distributed backend on \textbf{Microsoft Azure} integrating with \textbf{Microsoft CRM}, with event-driven processing, fault-tolerant retries over \textbf{SQL Server}, and parallelized batch ingestion that materially cut execution time. The same posture, careful service design backed by reliable cloud infrastructure, carries over to any production Java + AWS environment.

On testing and reliability: I'm comfortable navigating large codebases, debugging incidents through logs, metrics, and traces, and writing automated tests as part of the build rather than an afterthought.

Stack-wise: \textbf{Java}, \textbf{Python}, \textbf{TypeScript / JavaScript}, \textbf{SQL}, plus \textbf{AWS} (CDK, Lambda, API Gateway, DynamoDB, S3, IAM, CloudWatch, ECS Fargate, SQS, SNS), \textbf{Docker}, \textbf{Kubernetes}, and \textbf{CI/CD} pipelines. Comfortable with \textbf{Git} workflows and microservices design.

What pulls me toward Diverse Lynx is the operational seriousness of consulting-style delivery. Working across client engagements where the engineering has to be both production-quality and shippable on real timelines is the kind of accountability I want to keep building under.

I'd welcome the chance to discuss further. Thanks for considering my application.

\vspace{12pt}
Sincerely, \\
Ajay Venkatesh

\end{document}